\documentclass[a4paper,12pt,twoside]{memoir}

% Castellano
\usepackage[spanish,es-tabla]{babel}
\selectlanguage{spanish}
\usepackage[utf8]{inputenc}
\usepackage[T1]{fontenc}
\usepackage{lmodern} % scalable font
\usepackage{microtype}
\usepackage{placeins}

\RequirePackage{booktabs}
\RequirePackage[table]{xcolor}
\RequirePackage{xtab}
\RequirePackage{multirow}

% Links
\PassOptionsToPackage{hyphens}{url}\usepackage[colorlinks]{hyperref}
\hypersetup{
	allcolors = {red}
}

% Ecuaciones
\usepackage{amsmath}

% Rutas de fichero / paquete
\newcommand{\ruta}[1]{{\sffamily #1}}

% Párrafos
\nonzeroparskip

% Huérfanas y viudas
\widowpenalty100000
\clubpenalty100000

% Evitar solapes en el header
\nouppercaseheads

% Imagenes
\usepackage{graphicx}
\newcommand{\imagen}[2]{
	\begin{figure}[!h]
		\centering
		\includegraphics[width=0.9\textwidth]{#1}
		\caption{#2}\label{fig:#1}
	\end{figure}
	\FloatBarrier
}

\newcommand{\imagenflotante}[2]{
	\begin{figure}%[!h]
		\centering
		\includegraphics[width=0.9\textwidth]{#1}
		\caption{#2}\label{fig:#1}
	\end{figure}
}



% El comando \figura nos permite insertar figuras comodamente, y utilizando
% siempre el mismo formato. Los parametros son:
% 1 -> Porcentaje del ancho de página que ocupará la figura (de 0 a 1)
% 2 --> Fichero de la imagen
% 3 --> Texto a pie de imagen
% 4 --> Etiqueta (label) para referencias
% 5 --> Opciones que queramos pasarle al \includegraphics
% 6 --> Opciones de posicionamiento a pasarle a \begin{figure}
\newcommand{\figuraConPosicion}[6]{%
  \setlength{\anchoFloat}{#1\textwidth}%
  \addtolength{\anchoFloat}{-4\fboxsep}%
  \setlength{\anchoFigura}{\anchoFloat}%
  \begin{figure}[#6]
    \begin{center}%
      \Ovalbox{%
        \begin{minipage}{\anchoFloat}%
          \begin{center}%
            \includegraphics[width=\anchoFigura,#5]{#2}%
            \caption{#3}%
            \label{#4}%
          \end{center}%
        \end{minipage}
      }%
    \end{center}%
  \end{figure}%
}

%
% Comando para incluir imágenes en formato apaisado (sin marco).
\newcommand{\figuraApaisadaSinMarco}[5]{%
  \begin{figure}%
    \begin{center}%
    \includegraphics[angle=90,height=#1\textheight,#5]{#2}%
    \caption{#3}%
    \label{#4}%
    \end{center}%
  \end{figure}%
}
% Para las tablas
\newcommand{\otoprule}{\midrule [\heavyrulewidth]}
%
% Nuevo comando para tablas pequeñas (menos de una página).
\newcommand{\tablaSmall}[5]{%
 \begin{table}
  \begin{center}
   \rowcolors {2}{gray!35}{}
   \begin{tabular}{#2}
    \toprule
    #4
    \otoprule
    #5
    \bottomrule
   \end{tabular}
   \caption{#1}
   \label{tabla:#3}
  \end{center}
 \end{table}
}

%
%Para el float H de tablaSmallSinColores
\usepackage{float}

%
% Nuevo comando para tablas pequeñas (menos de una página).
\newcommand{\tablaSmallSinColores}[5]{%
 \begin{table}[H]
  \begin{center}
   \begin{tabular}{#2}
    \toprule
    #4
    \otoprule
    #5
    \bottomrule
   \end{tabular}
   \caption{#1}
   \label{tabla:#3}
  \end{center}
 \end{table}
}

\newcommand{\tablaApaisadaSmall}[5]{%
\begin{landscape}
  \begin{table}
   \begin{center}
    \rowcolors {2}{gray!35}{}
    \begin{tabular}{#2}
     \toprule
     #4
     \otoprule
     #5
     \bottomrule
    \end{tabular}
    \caption{#1}
    \label{tabla:#3}
   \end{center}
  \end{table}
\end{landscape}
}

%
% Nuevo comando para tablas grandes con cabecera y filas alternas coloreadas en gris.
\newcommand{\tabla}[6]{%
  \begin{center}
    \tablefirsthead{
      \toprule
      #5
      \otoprule
    }
    \tablehead{
      \multicolumn{#3}{l}{\small\sl continúa desde la página anterior}\\
      \toprule
      #5
      \otoprule
    }
    \tabletail{
      \hline
      \multicolumn{#3}{r}{\small\sl continúa en la página siguiente}\\
    }
    \tablelasttail{
      \hline
    }
    \bottomcaption{#1}
    \rowcolors {2}{gray!35}{}
    \begin{xtabular}{#2}
      #6
      \bottomrule
    \end{xtabular}
    \label{tabla:#4}
  \end{center}
}

%
% Nuevo comando para tablas grandes con cabecera.
\newcommand{\tablaSinColores}[6]{%
  \begin{center}
    \tablefirsthead{
      \toprule
      #5
      \otoprule
    }
    \tablehead{
      \multicolumn{#3}{l}{\small\sl continúa desde la página anterior}\\
      \toprule
      #5
      \otoprule
    }
    \tabletail{
      \hline
      \multicolumn{#3}{r}{\small\sl continúa en la página siguiente}\\
    }
    \tablelasttail{
      \hline
    }
    \bottomcaption{#1}
    \begin{xtabular}{#2}
      #6
      \bottomrule
    \end{xtabular}
    \label{tabla:#4}
  \end{center}
}

%
% Nuevo comando para tablas grandes sin cabecera.
\newcommand{\tablaSinCabecera}[5]{%
  \begin{center}
    \tablefirsthead{
      \toprule
    }
    \tablehead{
      \multicolumn{#3}{l}{\small\sl continúa desde la página anterior}\\
      \hline
    }
    \tabletail{
      \hline
      \multicolumn{#3}{r}{\small\sl continúa en la página siguiente}\\
    }
    \tablelasttail{
      \hline
    }
    \bottomcaption{#1}
  \begin{xtabular}{#2}
    #5
   \bottomrule
  \end{xtabular}
  \label{tabla:#4}
  \end{center}
}



\definecolor{cgoLight}{HTML}{EEEEEE}
\definecolor{cgoExtralight}{HTML}{FFFFFF}

%
% Nuevo comando para tablas grandes sin cabecera.
\newcommand{\tablaSinCabeceraConBandas}[5]{%
  \begin{center}
    \tablefirsthead{
      \toprule
    }
    \tablehead{
      \multicolumn{#3}{l}{\small\sl continúa desde la página anterior}\\
      \hline
    }
    \tabletail{
      \hline
      \multicolumn{#3}{r}{\small\sl continúa en la página siguiente}\\
    }
    \tablelasttail{
      \hline
    }
    \bottomcaption{#1}
    \rowcolors[]{1}{cgoExtralight}{cgoLight}

  \begin{xtabular}{#2}
    #5
   \bottomrule
  \end{xtabular}
  \label{tabla:#4}
  \end{center}
}




\graphicspath{ {./img/} }

% Capítulos
\chapterstyle{bianchi}
\newcommand{\capitulo}[2]{
	\setcounter{chapter}{#1}
	\setcounter{section}{0}
	\setcounter{figure}{0}
	\setcounter{table}{0}
	\chapter*{#2}
	\addcontentsline{toc}{chapter}{#2}
	\markboth{#2}{#2}
}

% Apéndices
\renewcommand{\appendixname}{Apéndice}
\renewcommand*\cftappendixname{\appendixname}

\newcommand{\apendice}[1]{
	%\renewcommand{\thechapter}{A}
	\chapter{#1}
}

\renewcommand*\cftappendixname{\appendixname\ }

% Formato de portada
\makeatletter
\usepackage{xcolor}
\newcommand{\tutor}[1]{\def\@tutor{#1}}
\newcommand{\course}[1]{\def\@course{#1}}
\definecolor{cpardoBox}{HTML}{E6E6FF}
\def\maketitle{
  \null
  \thispagestyle{empty}
  % Cabecera ----------------
\noindent\includegraphics[width=\textwidth]{cabecera}\vspace{1cm}%
  \vfill
  % Título proyecto y escudo informática ----------------
  \colorbox{cpardoBox}{%
    \begin{minipage}{.8\textwidth}
      \vspace{.5cm}\Large
      \begin{center}
      \textbf{TFG del Grado en Ingeniería Informática}\vspace{.6cm}\\
      \textbf{\LARGE\@title{}}
      \end{center}
      \vspace{.2cm}
    \end{minipage}

  }%
  \hfill\begin{minipage}{.20\textwidth}
    \includegraphics[width=\textwidth]{escudoInfor}
  \end{minipage}
  \vfill
  % Datos de alumno, curso y tutores ------------------
  \begin{center}%
  {%
    \noindent\LARGE
    Presentado por \@author{}\\ 
    en Universidad de Burgos --- \@date{}\\
    Tutor: \@tutor{}\\
  }%
  \end{center}%
  \null
  \cleardoublepage
  }
\makeatother


% Datos de portada
\title{título del TFG \\Documentación Técnica}
\author{nombre alumno}
\tutor{nombre tutor}
\date{\today}

\begin{document}

\maketitle



\cleardoublepage



%%%%%%%%%%%%%%%%%%%%%%%%%%%%%%%%%%%%%%%%%%%%%%%%%%%%%%%%%%%%%%%%%%%%%%%%%%%%%%%%%%%%%%%%



\frontmatter


\clearpage

% Indices
\tableofcontents

\clearpage

\listoffigures

\clearpage

\listoftables

\clearpage

\mainmatter

\appendix

\apendice{Plan de Proyecto Software}

\section{Introducción}

En este apartado se procede a explicar la planificación temporal del proyecto, detallando cada uno de los \textit{sprints} desde el comienzo del proyecto hasta el final del mismo. También se aborda un estudio de viabilidad del proyecto, tanto económica como legal.

\section{Planificación temporal}

Como ya se ha explicado en la sección de técnica metodológica, la distribución temporal del proyecto se ha realizado mediante \textit{Sprints} de aproximadamente dos semanas de duración.
Durante los primeros seis \textit{Sprints} se ha usado la herramienta de gestión de tareas \textit{Microsoft Planner}

\subsection{Sprint 1: Primeros pasos (23/11/2023 - 20/12/2023)}
Durante la primera reunión se barajaron diferentes propuestas de proyecto, pero hubo una que me despertó especial interés y me acabé decantando por ella. \textit{Open Data in Smart Cities}.
Se establecieron las primeras tareas, como investigar acerca de \textit{FIWARE}, descargar la plantilla de memoria, pensar en las diferentes herramientas para la organización del proyecto (Repositorio, planificador de tareas, control de versiones...). 

Durante el Sprint, se terminaron por decidir dichas herramientas, se creó el repositorio y se incorporaron en el planificador de tareas las primeras tareas para empezar con la investigación y profundización en el ámbito del proyecto.
También se comenzó con la investigación de \textit{Flask}, el \textit{framework} usado para el desarrollo de la aplicación en \textit{Python}. Como finalmente se aplazó la reunión programada para esas dos semanas, se comenzó con una pequeña demo de tratamiento de datos con \textit{Flask}.

\imagen{Sprint 1}{Sprint 1 - Tareas}

\subsection{Sprint 2: Investigación \textit{FIWARE} y desarrollo de la aplicación (20/12/2023 - 17/01/2024)}

Durante este Sprint se continuó con la investigación acerca de \textit{FIWARE}, y se trabajó en un menú para el usuario y el tratamiento de datos.
También se comenzó con una demo paralela con interfaces gráficas, para decidir si finalmente se usaba \textit{Flask} o se desarrollaba la aplicación con interfaces gráficas.

\imagen{Sprint 2}{Sprint 2 - Tareas}

\subsection{Sprint 3: Despliegue del servidor OCB (17/01/2024 - 31/01/2024)}

En esta etapa se investigó acerca del funcionamiento de \textit{Docker} y se comenzó con el despliegue del \textit{Orion Context Broker}.
También se creó un catalogo de datos anexado en la propia tarea del \textit{Planner}, el cual contiene los diferentes enlaces y los formatos que contenían. Aquí se empezó a replantear el uso de del modelo de información \textit{NGSI-LD}, debido a la falta de conjuntos con ese modelo.
También se desarrollo un \textit{Script} que transformaba ficheros en formato \textit{JSON} al modelo \textit{NGSI-LD}, dicho \textit{script} sería descartado posteriormente por su escasa eficacia y funcionalidad.

\imagen{Sprint 3}{Sprint 3 - Tareas}

\subsection{Sprint 4: Despliegue total y testeo del \textit{OCB} (31/01/2024 - 14/02/2024)}

Durante esta etapa se siguió con el desarrollo del servidor del \textit{Orion Context Broker} con el apoyo de la herramienta \textit{Postman}, con la cuál se podía ir comprobando que la base de datos está guardando bien las entidades y es capaz de devolverlas.

\imagen{Sprint 4}{Sprint 4 - Tareas}

\subsection{Sprint 5: Simulador (14/02/2024 - 14/03/2024)}

En esta etapa del proyecto afrontamos el problema de la falta de datos con el modelo \textit{NGSI-LD}, desarrollando un \textit{script} que simula un conjunto de datos con este modelo.
Este simulador representaba un conjunto de \textit{parkings}, cuyos sensores publican cada 30 segundos el estado actual de las entidades (\textit{parkings}) en la base de datos, para ser consultadas desde la aplicación.

Al no tener una caché el \textit{OCB}, no se podía guardar el estado anterior de las entidades, por lo que no se tenía un histórico de las mismas, esta funcionalidad aportaría grandes funcionalidades como análisis de patrones o predecir el estado de las entidades. Por lo tanto, durante esta etapa se investigó acerca de alguna opción que permitiese guardar históricos de las entidades.

Existe una herramienta que forma parte de la plataforma \textit{FIWARE}, llamada \textit{Cygnus}, la cual permite guardar las actualizaciones de las entidades. Esta herramienta se intentó implementar para este caso y por complejidad y poca funcionalidad (ya que es solo un simulador y se puede ir guardando un histórico antes de publicar los datos en el servidor, agilizando así la aplicación), no se terminó de implantar. Pero puede ser útil para futuras versiones.

También durante Sprint se desarrolló una demo de una plantilla de \textit{PowerBi} para dichos datos, para intentar integrarla en aplicación. El inconveniente que se encontró es que \textit{PowerBi} es un software de pago, por lo tanto, no se pudo finalmente integrar en la aplicación.


\imagen{Sprint 5}{Sprint 5 - Tareas}

\subsection{Sprint 6: Nuevas funcionalidades en la aplicación (14/03/2024 - 24/04/2024)}

Durante esta etapa la aplicación ya estaba cogiendo forma, por lo que se comenzó con los cambios estéticos, para ofrecer al usuario una mayor simpleza y pueda tener una experiencia más amigable con la aplicación.
También se añadieron algunas funcionalidades como la de los registros en la aplicación, o alguna plantilla más visual con la librería \textit{Dash}.

También se investigó e incorporó la manera de autentificar a los usuarios con un \textit{login}.

\imagen{Sprint 6}{Sprint 6 - Tareas}

\subsection{Sprint 7: (24/04/2024 - 15/05/2024)}

A partir de aquí se llevó a cabo un cambio de herramienta para la gestión de tareas. Hasta entonces se había estado usando \textit{Microsoft Planner}, pero surgió una alternativa, la cual es mucho más práctica porque permitía enlazar las tareas con los \textit{commits} y tener un mejor control de las versiones y de las tareas. Esta alternativa es ofrecida por \textit{GitHub}, nuestro repositorio y herramienta de control de versiones. Mediante los \textit{issues}, los cuales sustituyen a las tareas en el \textit{Planner} y los \textit{milestones}, los cuales representan cada uno de los \textit{sprints}, se podía lograr una mejor integración de las herramientas y una simplificación del trabajo muy destacable.

En esta etapa se desarrolló una pequeña aplicación para la gestión del servidor \textit{OCB}. Como se muestra en la imagen \ref{fig:GestorOCB}, esta aplicación permite:
\begin{itemize}
    \item \textbf{Instalar o desinstalar el servidor:} permite el la instalación o desinstalación del \textit{docker}.
    \item \textbf{Iniciar o parar el servidor:} permite iniciar o parar la ejecución del \textit{docker} de \textit{orion} y el de la \textit{MongoDB}, y eliminar su conexión.
    \item \textbf{Mostrar o eliminar:} permite mostrar o eliminar las entidades almacenadas en el servidor.
    \item \textbf{Iniciar el simulador:} inicia el simulador del modelo \textit{NGSI-LD}.
\end{itemize}

\imagen{GestorOCB}{Menú para la gestión del \textit{OCB}}
También se incorporaron nuevas funcionalidades en la aplicación como la posibilidad de añadir nuevos conjuntos de distintos formatos y una plantilla genérica para estos casos.
Por último, se trabajó en uno de los objetivos clave de la aplicación, conseguir que el usuario tenga una experiencia cómoda y sencilla con la aplicación. Para ello se investigó y se integró\textit{Bootstrap}, y también, se desarrolló una guía de usuario con \textit{GitBook}.

\imagen{Sprint 7}{Sprint 7 - Tareas}

\subsection{Sprint 8: Organización de la aplicación (15/05/2024 - 31/05/2024)}

Durante esta etapa se llevó a cabo un proceso de organización del desarrollo de la aplicación, mediante el cuál se elaboraron los diagramas de ventanas, casos de uso, interacción. Pudiendo así, enfocar la aplicación en unos requisitos determinados y simplificar algunas funcionalidades.
Durante el proceso se añadieron funcionalidades como añadir mapas en algunas vistas, diccionarios de traducciones e interpretación de datos y mejoras visuales para el usuario.

\imagen{Sprint 8}{Sprint 8 - Tareas}

\subsection{Sprint 9: Últimas funcionalidades (31/05/2024 - 14/06/2024)}

En esta etapa del proyecto se comenzó a profundizar un poco más en la redacción de la memoria y se añadieron las últimas funcionalidades en la aplicación como por ejemplo:
\begin{itemize}
    \item \textbf{Funcionalidades de administrador:} eliminar usuarios y cambiar sus roles.
    \item \textbf{Favoritos:} posibilidad de marcar conjuntos como favoritos.
    \item \textbf{\textit{Tooltip} dinámico:} en los mapas de los conjuntos, el usuario puede seleccionar que parámetros visualizar.
\end{itemize}

También se investigaron otras funcionalidades que por complejidad o falta de tiempo se terminaron descartando, como añadir filtros para los parámetros de los conjuntos o una barra de tiempo.
Otros punto a destacar en este \textit{sprint} fueron los cambios estéticos en la aplicación.

\imagen{Sprint 9}{Sprint 9 - Tareas}

\subsection{Sprint 10: Primer tramo de la recta final (14/06/2024 - 25/06/2024)}

Estando próxima la entrega, durante este \textit{sprint} se avanzó con los distintos apartados de la memoria y de los anexos.
En cuanto a la aplicación, se incorporaron mejoras visuales, como la homogeneización de la aplicación aplicando un estilo común, ayudado de \textit{Bootstrap}. Se procedió también a la revisión y mejora de algunas funcionalidades, como la gestión de registros. Y por último se incorporó un historial de errores en el cual el usuario podía dejar constancia de los distintos errores que pudieran ocurrir durante la aplicación.

\imagen{Sprint 10}{Sprint 10 - Tareas}

\subsection{Sprint 11: Últimos retoques y anexos (25/06/2024 - 03/07/2024)}

AÑADIR

\section{Estudio de viabilidad}

\subsection{Viabilidad económica}

En este apartado se va a analizar los costes y beneficios del desarrollo de este proyecto.

\subsubsection{Costes}
 
Este proyecto ha llevado más de 500 horas de trabajo:\begin{itemize}
    \item \textbf{480 horas} de trabajo por parte del alumno.
    \item \textbf{20 horas} de reuniones con los tutores.
\end{itemize}  

Para extrapolar estas cifras a un proyecto real, se identificarán a los tutores como \textit{Project Mangers} y al alumno como el resto de personal que formaría parte de un equipo de desarrollo (analista, diseñador, programador...).

Un \textit{Project Manger} en España cobra de media 40.000€ anuales\footnote{\url{https://www.glassdoor.es/Sueldos/project-manager-sueldo-SRCH_KO0,15.htm}}, teniendo en cuenta que la jornada laboral en España normalmente es de 8 horas al día durante 5 días a la semana, el sueldo por hora sería de 19,23€\footnote{Conversión realizada con \url{https://convertir-sueldo-hora-ano.appspot.com}}.
Si multiplicamos las 20 horas totales de cada \textit{Project Manger} por el salario por hora medio en España de un \textit{Project Manger} por los dos \textit{Project Mangers} del proyecto, obtenemos un total de 769,2€

$$20 \, \text{horas} \times 19,23 \, \text{€/hora} \times 2 \, \text{personas} = \text{769,2€}$$

El sueldo medio de un desarrollador en España es de 20.000€ anuales\footnote{\url{https://www.hackaboss.com/blog/cuanto-gana-programador-salario-espana}}, volviendo a hacer la conversión a euros la hora, sería 9,62€ por hora\footnote{Conversión realizada con \url{https://convertir-sueldo-hora-ano.appspot.com}}.
Si multiplicamos las 480 horas totales del proyecto por el salario por hora medio en España de un desarrollador, obtenemos un total de 4.617,60€

$$480 \, \text{horas} \times 9,62 \, \text{€/hora} = \text{4.617,60€}$$

En total el coste de personal sería de 5.386,8€
$$769,2 \, \text{€} + 4.617,60 \, \text{€} = \text{5.386,8€}$$

El equipo utilizado es un ordenador Lenovo Legion Y520 con un procesador Intel® Core™ i7-7700HQ de 3.8GHz con 8 GB de RAM, valorado en 799,99€\footnote{\url{https://www.mediamarkt.es/es/product/_portátil-gaming-lenovo-legion-y520-15ikbn-15-6-fhd-intel®-core™-i7-7700hq-8gb-1tb-gtx1050-fdos-1401315.html}}.

En cuanto a \textit{software}, todas las herramientas son gratuitas. Solo se ha usado de pago la herramienta \LaTeX{}, que aunque es gratuita, el último mes se ha decidido probar por 8€ su versión Pro, para poder mejorar la conectividad con \textit{GitHub} durante la redacción de la memoria.

En total el coste del \textit{hardware} y \textit{software} sería de 5.386,8€
$$799,99 \, \text{€} + 8 \, \text{€} = \text{807,99€}$$

Si se suma el coste de personal con el de \textit{hardware} y \textit{software}, sale un total de 6.194,79€
$$5.386,8 \, \text{€} + 807,99 \, \text{€} = \text{6.194,79€}$$

\subsubsection{Beneficios}

El proyecto está enfocado a acercar el acceso a los datos en tiempo real a la población de una manera no lucrativa, es por ello que la aplicación es totalmente gratuita.
Una opción que se podría incluir para cubrir los gastos del mantenimiento de la aplicación puede ser incluir un apartado de donaciones, para que los usuarios que quieran contribuir puedan hacerlo de manera altruista.

Por otro lado, si fuese necesario se podría añadir anuncios de manera que generen algunos beneficios para el mantenimiento de la aplicación, pero en ningún caso se plantean fines lucrativos con la aplicación.

\subsection{Viabilidad legal}

Al ser una aplicación pensada para ser usada en cualquier ciudad del mundo, es muy difícil tratar de realizar un estudio legal sobre el proyecto, ya que, cada país tiene su propias leyes de protección de datos o acceso a datos públicos. Por lo tanto, se va a tomar como referencia las leyes de la legislación española, lugar donde se realiza este proyecto.
Los principales puntos a tratar para mantener la aplicación dentro del marco legal españon son los siguientes:
\begin{itemize}
    \item \textbf{Derecho de Acceso a la Información Pública:} En España, el derecho de acceso a la información pública está regulado por la Ley 19/2013, de 9 de diciembre, de transparencia, acceso a la información pública y buen gobierno \cite{ley-transparencia}.

    \item \textbf{Licencias de Datos Abiertos:} Los datos abiertos proporcionados por instituciones públicas en España suelen estar bajo la licencia Open Government License - España \cite{ogl-es}, la cual permite su uso, redistribución y modificación bajo ciertas condiciones.

    \item \textbf{Protección de Datos Personales:} El tratamiento de datos personales en España está regulado por la Ley Orgánica 3/2018, de 5 de diciembre, de Protección de Datos Personales y garantía de los derechos digitales \cite{lopdgdd}. Es crucial cumplir con los principios de protección de datos, obtener consentimiento cuando sea necesario y garantizar la seguridad de los datos personales.

    \item \textbf{Derechos de Autor:} En España, los datos en sí no están protegidos por derechos de autor, pero las bases de datos pueden tener protección si cumplen con los requisitos establecidos en la Ley de Propiedad Intelectual \cite{lpi}.

    \item \textbf{Normativas Aplicables a Aplicaciones Web:} Además de las leyes específicas mencionadas, la aplicación debe cumplir con la normativa general aplicable a las aplicaciones web en términos de comercio electrónico, protección de consumidores y usuarios, entre otros.
\end{itemize}

Es importante considerar las licencias de las librerías y otros componentes de software utilizados en la aplicación. En la tabla \ref{tabla-licencias} se detallan las librerías utilizadas y sus respectivas licencias:

\begin{table}
	\scalebox{0.80}{
		\begin{tabular}{@{}p{14em} p{6em} p{20em}@{}}
			\toprule
			\textbf{Librería} & \textbf{Versión} & \textbf{Licencia}\\
			\midrule
            \texttt{\textbf{blinker}} & \texttt{1.8.2} & MIT License \\
            \texttt{\textbf{certifi}} & \texttt{2024.7.4} & Mozilla Public License 2.0 (MPL-2.0) \\
            \texttt{\textbf{charset-normalizer}} & \texttt{3.3.2} & MIT License \\
            \texttt{\textbf{click}} & \texttt{8.1.7} & BSD-3-Clause License \\
            \texttt{\textbf{colorama}} & \texttt{0.4.6} & BSD-3-Clause License \\
            \texttt{\textbf{dash}} & \texttt{2.17.1} & MIT License \\
            \texttt{\textbf{dash-core-components}} & \texttt{2.0.0} & MIT License \\
            \texttt{\textbf{dash-html-components}} & \texttt{2.0.0} & MIT License \\
            \texttt{\textbf{dash-table}} & \texttt{5.0.0} & MIT License \\
            \texttt{\textbf{Flask}} & \texttt{3.0.3} & BSD-3-Clause License \\
            \texttt{\textbf{Flask-Login}} & \texttt{0.6.3} & MIT License \\
            \texttt{\textbf{Flask-MySQLdb}} & \texttt{2.0.0} & MIT License \\
            \texttt{\textbf{Flask-WTF}} & \texttt{1.2.1} & BSD-3-Clause License \\
            \texttt{\textbf{idna}} & \texttt{3.7} & BSD-3-Clause License \\
            \texttt{\textbf{importlib\_metadata}} & \texttt{8.0.0} & Apache License 2.0 \\
            \texttt{\textbf{itsdangerous}} & \texttt{2.2.0} & BSD-3-Clause License \\
            \texttt{\textbf{Jinja2}} & \texttt{3.1.4} & BSD-3-Clause License \\
            \texttt{\textbf{MarkupSafe}} & \texttt{2.1.5} & BSD-3-Clause License \\
            \texttt{\textbf{mysql-connector-python}} & \texttt{9.0.0} & GPL-2.0 License \\
            \texttt{\textbf{mysqlclient}} & \texttt{2.2.4} & GPL-2.0 License \\
            \texttt{\textbf{nest-asyncio}} & \texttt{1.6.0} & BSD-3-Clause License \\
            \texttt{\textbf{numpy}} & \texttt{2.0.0} & BSD-3-Clause License \\
            \texttt{\textbf{packaging}} & \texttt{24.1} & Apache License 2.0 \\
            \texttt{\textbf{pandas}} & \texttt{2.2.2} & BSD-3-Clause License \\
            \texttt{\textbf{plotly}} & \texttt{5.22.0} & MIT License \\
            \texttt{\textbf{python-dateutil}} & \texttt{2.9.0.post0} & BSD License \\
            \texttt{\textbf{pytz}} & \texttt{2024.1} & MIT License \\
            \texttt{\textbf{requests}} & \texttt{2.32.3} & Apache License 2.0 \\
            \texttt{\textbf{retrying}} & \texttt{1.3.4} & Apache License 2.0 \\
            \texttt{\textbf{setuptools}} & \texttt{70.2.0} & MIT License \\
            \texttt{\textbf{six}} & \texttt{1.16.0} & MIT License \\
            \texttt{\textbf{tenacity}} & \texttt{8.4.2} & MIT License \\
            \texttt{\textbf{typing\_extensions}} & \texttt{4.12.2} & MIT License \\
            \texttt{\textbf{tzdata}} & \texttt{2024.1} & MIT License \\
            \texttt{\textbf{urllib3}} & \texttt{2.2.2} & MIT License \\
            \texttt{\textbf{Werkzeug}} & \texttt{3.0.3} & BSD-3-Clause License \\
            \texttt{\textbf{WTForms}} & \texttt{3.1.2} & BSD-3-Clause License \\
            \texttt{\textbf{xmltodict}} & \texttt{0.13.0} & MIT License \\
            \texttt{\textbf{zipp}} & \texttt{3.19.2} & MIT License \\
			\bottomrule
		\end{tabular}
	}
	\caption{Tabla de licencias}
    \label{tabla-licencias}
\end{table}

Entre las licencias mencionadas para las librerías, la más restrictiva es la GPL-2.0 License (General Public License versión 2.0). Esta licencia impone fuertes requisitos de copyleft, lo que significa que cualquier software derivado de una obra licenciada bajo la GPL-2.0 también debe ser distribuido bajo la misma licencia. Esto no se aplica a las licencias más permisivas como MIT, BSD-3-Clause, o Apache License 2.0, que permiten una mayor flexibilidad en el uso y redistribución del software.
Por lo tanto la licencia escogida para la distribución del software es la licencia  GPL-2.0 License (General Public License versión 2.0)\footnote{\url{https://www.gnu.org/licenses/old-licenses/gpl-2.0.html#SEC1}}.



\apendice{Especificación de Requisitos}

\section{Introducción}

En este apartado se van a explicar los distintos requisitos de la aplicación, marcados principalmente por los objetivos establecidos \ref{objetivos-generales}.
Estos requisitos son la base de la aplicación, ya que actúan como guía después de la etapa de análisis para el correcto desarrollo de la aplicación con un enfoque claro y unos objetivos concretos.

El anexo está estructurado de la siguiente forma:
\begin{itemize}
    \item Objetivos generales: contiene los objetivos generales del proyecto.
    \item Catálogo de requisitos: se detallan los requisitos funcionales y no funcionales de la aplicación 
    \item Especificación de requisitos: se muestra una visión general de los casos de uso del sistema y posteriormente se detallan uno a uno.
\end{itemize}

\section{Objetivos generales}\label{objetivos-generales}

Los objetivos generales del proyecto son los siguientes:

\begin{itemize}
    \item Desarrollar una aplicación que permita a cualquier usuario el acceso a los distintos conjuntos de datos, que los servicios de \textit{IT} hayan publicado en abierto, de manera visual, personalizada y precisa.
    \item Ofrecer al usuario una experiencia intuitiva y cómoda con la aplicación.
    \item Permitir al usuario la posibilidad de guardar un histórico de los conjuntos que desee para un posterior análisis de los mismos.
    \item Contribuir con el crecimiento de la plataforma \textit{FIWARE}.

\end{itemize}

\section{Catálogo de requisitos}\label{catalogo-de-requisitos}

A continuación se detallan los requisitos funcionales y no funcionales de la aplicación.

\subsection{\textbf{Requisitos funcionales}}

\begin{itemize}
    \item \textbf{RF-1 Inicio de sesión:} El Sistema debe permitir al usuario iniciar sesión introduciendo su usuario y su contraseña.
    \begin{itemize}
        \item \textbf{RF-1.1 Privacidad de contraseña:} Debe poder ocultarse o mostrarse la contraseña.
        \item \textbf{RF-1.2 Contraseña incorrecta:} Debe notificarse al usuario si la contraseña no es correcta.
        \item \textbf{RF-1.3 Usuario inexistente:} Debe notificarse al usuario si el nombre introducido no se encuentra en el Sistema.
    \end{itemize}

    \item \textbf{RF-2 Crear una cuenta:} El Sistema debe permitir al usuario crear una cuenta introduciendo un nombre de usuario, su nombre completo y una contraseña.
    \begin{itemize}
        \item \textbf{RF-2.1 Verificación de usuario:} Debe notificarse al usuario si ya existe un usuario con ese nombre de usuario en el Sistema.
        \item \textbf{RF-2.2 Privacidad de incorrecta:} Debe poder ocultarse o mostrarse la contraseña.
        \item \textbf{RF-2.3 Verificación de contraseña:} Debe comprobarse que las contraseñas coincidad y en caso contrario notificarlo al usuario.
    \end{itemize}

    \item \textbf{RF-3 Elegir conjuntos favoritos:} El usuario debe poder marcar o desmarcar cualquier conjunto como favorito desde la vista de dicho conjunto.
    \begin{itemize}
        \item \textbf{RF-3.1 Acceso rápido:} Los conjuntos marcados como favoritos aparecerán en la pantalla de inicio con el objetivo de poder acceder a ellos de manera rápida.
    \end{itemize}

    \item \textbf{RF-4 Gestión conjuntos:} El usuario debe poder añadir conjuntos de datos en distintos formatos a la aplicación para que el resto de usuarios puedan consultarlos.
    \begin{itemize}
        \item \textbf{RF-4.1 Campos a introducir:} El usuario debe introducir obligatoriamente los campos de ciudad, conjunto, formato y enlace, y opcionalmente la periodicidad con la que se actualiza dicho conjunto.
        \item \textbf{RF-4.2 Eliminación de conjuntos:} El administrador debe poder eliminar los conjuntos que considere pertinentes.
    \end{itemize}

    \item \textbf{RF-5 Gestión de registros:} El Sistema debe ofrecer al usuario la posibilidad de guardar un histórico de los conjuntos con una periodicidad indicada por el usuario.
    \begin{itemize}
        \item \textbf{RF-5.1 Registros repetidos:} La periodicidad no puede ser menor a la indicada en origen del conjunto de datos.
        \item \textbf{RF-5.2 Descarga de registros:} El usuario debe poder descargar un conjunto de datos de los que esté guardando.
        \item \textbf{RF-5.3 Eliminación de registros:} El usuario debe poder eliminar el registro de un conjunto de datos.
    \end{itemize}

    \item \textbf{RF-6 Gestión de usuarios:} los usuarios con el rol de administrador deben poder gestionar el resto de usuarios.
    \begin{itemize}
        \item \textbf{RF-6.1 Cambio de rol:} los administradores deben poder cambiar el rol de cualquier usuario.
        \item \textbf{RF-6.2 Eliminar usuarios:} los administradores deben poder eliminar cualquier cuenta, menos la suya.
    \end{itemize}

    \item \textbf{RF-7 Traducción de campos:} los campos de los conjuntos de datos deben poder ser traducidos para una experiencia más cómoda del usuario.
    \begin{itemize}
        \item \textbf{RF-7.1 Gestionar traducciones:} los administradores deben poder añadir o eliminar traducciones de la aplicación, las cuales se aplican a todos los conjuntos.
    \end{itemize}

    \item \textbf{RF-8 Visualizar conjuntos de datos en tiempo real:} el usuario debe poder visualizar un conjunto de datos de una forma funcional y cómoda.
    \begin{itemize}
        \item \textbf{RF-8.1 Ordenación de entidades:} los usuarios deben poder ordenar alfabéticamente ascendente y descendente las distintas entidades mostradas en función de los parámetros de las mismas.
        \item \textbf{RF-8.2 Modificación del \textit{tooltip} del mapa:} los usuarios deben poder indicar que parámetros desean mostrar u ocultar en el mapa.
    \end{itemize}

    \item \textbf{RF-9 Reporte de errores:} el usuario debe poder reportar los errores que le impidan tener un uso adecuado de la aplicación.

    \item \textbf{RF-10 Mostrar ayuda:} el usuario debe poder consultar una pantalla de ayuda con instrucciones de ayuda y contacto con los administradores.
    
\end{itemize}

\subsection{\textbf{Requisitos no funcionales}}

\begin{itemize}
\tightlist
\item
    \textbf{RNF-1 Usabilidad:} la aplicación debe ofrecer una experiencia cómoda para el usuario, siendo lo más intuitiva y fácil de entender posible para el usuario.
\item
    \textbf{RNF-2 Compatibilidad:} la aplicación debe ser compatible con cualquier navegador en cualquier dispositivo.
\item
    \textbf{RNF-3 Disponibilidad:} la aplicación debe estar operativa en todo momento para los usar.
\item
    \textbf{RNF-4 Mantenibilidad y Escalabilidad}: la aplicación debe poder sen mantenible y fácilmente modificada por un desarrollador para aplicar futuras mejoras.
\item
    \textbf{RNF-5 Rendimiento}: la navegación por la aplicación debe ser ágil y los tiempos de carga de los conjuntos de datos deben ser muy reducidos.
\item
    \textbf{RNF-6 Persistencia}: los datos manejados relacionados a los usuarios deben de ser almacenados correctamente y de forma segura.
\item
    \textbf{RNF-7 Robustez:} la aplicación debe ser robusta y tener una respuesta y control a todos los errores para garantizar el correcto funcionamiento de la misma de forma continuada
\end{itemize}

\section{Especificación de requisitos}

Una vez detallados los requisitos funcionales de la aplicación, se detallan a continuación los casos de uso que proporcionan un contexto detallado a los requisitos del apartado anterior. \ref{catalogo-de-requisitos}

\subsection{Actores}
Los actores definidos en este sistema son:
\begin{itemize}
\tightlist
\item
    \textbf{Usuario:} cualquier usuario que se registre en la aplicación.
\item
    \textbf{Administrador:} usuario registrado en la aplicación con permisos de administrador, este rol solo puede ser asignado o denegado por otro administrador.
\end{itemize}

\subsection{Diagrama de casos de uso}
\imagen{Diagrama casos de uso}{Diagrama de casos de uso}

\subsection{Casos de uso}
% Caso de Uso 1 -> Crear una cuenta.
\begin{table}[p]
	\centering
	\begin{tabularx}{\linewidth}{ p{0.21\columnwidth} p{0.71\columnwidth} }
		\toprule
		\textbf{CU-1}    & \textbf{Crear una cuenta}\\
		\toprule
		\textbf{Versión}              & 1.0    \\
		\textbf{Autor}                & {Alejandro de Castro} \\
		\textbf{Requisitos asociados} & RF-2, RF-xx \\
		\textbf{Descripción}          & La descripción del CU \\
		\textbf{Precondición}         & Precondiciones (podría haber más de una) \\
		\textbf{Acciones}             &
		\begin{enumerate}
			\def\labelenumi{\arabic{enumi}.}
			\tightlist
			\item Pasos del CU
			\item Pasos del CU (añadir tantos como sean necesarios)
		\end{enumerate}\\
		\textbf{Postcondición}        & Postcondiciones (podría haber más de una) \\
		\textbf{Excepciones}          & Excepciones \\
		\textbf{Importancia}          & Alta o Media o Baja... \\
		\bottomrule
	\end{tabularx}
	\caption{CU-1 Nombre del caso de uso.}
\end{table}


% Caso de Uso 1 -> Crear cuenta.
\begin{table}[p]
	\centering
	\begin{tabularx}{\linewidth}{ p{0.21\columnwidth} p{0.71\columnwidth} }
		\toprule
		\textbf{CU-1}    & \textbf{Crear cuenta}\\
		\toprule
		\textbf{Versión}              & 1.0    \\
		\textbf{Autor}                & {Alejandro de Castro} \\
		\textbf{Requisitos asociados} & RF-2, RF-2.1, RF-2.2, RF-2.3 \\
		\textbf{Descripción}          & Permite al usuario crear una cuenta \\
		\textbf{Precondición}         & El usuario debe no haber iniciado sesión \\
		\textbf{Acciones}             &
		\begin{enumerate}
			\def\labelenumi{\arabic{enumi}.}
			\tightlist
			\item Rellenar los  campos de nombre de usuario, nombre completo, contraseña y confirmar contraseña.
			\item Pulsar en el botón de registrar.
		\end{enumerate}\\
		\textbf{Postcondición}        & Los datos son guardados en la base de datos y el usuario es creado, posteriormente es redirigido a la página de inicio de sesión. \\
		\textbf{Excepciones}          & En caso de no coincidir los campos de contraseñas o existir un usuario con ese nombre de usuario en la base de datos se notificará al usuario. \\
		\textbf{Importancia}          & Alta \\
		\bottomrule
	\end{tabularx}
	\caption{CU-1 Crear cuenta.}
\end{table}

% Caso de Uso 2 -> Iniciar sesión.
\begin{table}[p]
	\centering
	\begin{tabularx}{\linewidth}{ p{0.21\columnwidth} p{0.71\columnwidth} }
		\toprule
		\textbf{CU-2}    & \textbf{Iniciar sesión}\\
		\toprule
		\textbf{Versión}              & 1.0    \\
		\textbf{Autor}                & {Alejandro de Castro} \\
		\textbf{Requisitos asociados} & RF-1, RF-1.1, RF-1.2, RF-1.3 \\
		\textbf{Descripción}          & Permite al usuario iniciar sesión. \\
		\textbf{Precondición}         & \begin{itemize}
		    \item El usuario debe no haber iniciado sesión previamente. 
            \item El usuario debe tener una cuenta en el Sistema.
		\end{itemize} \\
		\textbf{Acciones}             &
		\begin{enumerate}
			\def\labelenumi{\arabic{enumi}.}
			\tightlist
			\item Rellenar los  campos de nombre de usuario y contraseña.
			\item Pulsar en el botón de iniciar sesión.
		\end{enumerate}\\
		\textbf{Postcondición}        & El usuario es redirigido a la página principal de la aplicación. \\
		\textbf{Excepciones}          & \begin{itemize}
		    \item Si el nombre de usuario no se encuentra en la base de datos se le notificará al usuario con un mensaje de error. 
            \item Si la contraseña es incorrecta se le notificará al usuario con un mensaje de error. 
		\end{itemize} \\
		\textbf{Importancia}          & Alta \\
		\bottomrule
	\end{tabularx}
	\caption{CU-2 Iniciar sesión.}
\end{table}

% Caso de Uso 3 -> Visualizar conjunto.
\begin{table}[p]
	\centering
	\begin{tabularx}{\linewidth}{ p{0.21\columnwidth} p{0.71\columnwidth} }
		\toprule
		\textbf{CU-3}    & \textbf{Visualizar conjunto}\\
		\toprule
		\textbf{Versión}              & 1.0    \\
		\textbf{Autor}                & {Alejandro de Castro} \\
		\textbf{Requisitos asociados} & RF-8 \\
		\textbf{Descripción}          & Permite al usuario visualizar un conjunto de datos. \\
		\textbf{Precondición}         & El usuario debe haber iniciado sesión. \\
		\textbf{Acciones}             &
		\begin{enumerate}
			\def\labelenumi{\arabic{enumi}.}
			\tightlist
			\item Seleccionar un conjunto de datos en la pantalla principal o en la página de conjuntos.
			\item Pulsar en el botón de consultar.
		\end{enumerate}\\
		\textbf{Postcondición}        & El usuario es redirigido a la página de visualización de conjuntos. \\
		\textbf{Excepciones}          & Si el conjunto de datos no está accesible se le redirigirá a la ventana de gestión de errores. \\
		\textbf{Importancia}          & Alta \\
		\bottomrule
	\end{tabularx}
	\caption{CU-3 Visualizar conjunto.}
\end{table}

% Caso de Uso 3.1 -> Personalizar conjunto.
\begin{table}[p]
	\centering
	\begin{tabularx}{\linewidth}{ p{0.21\columnwidth} p{0.71\columnwidth} }
		\toprule
		\textbf{CU-3.1}    & \textbf{Personalizar conjunto}\\
		\toprule
		\textbf{Versión}              & 1.0    \\
		\textbf{Autor}                & {Alejandro de Castro} \\
		\textbf{Requisitos asociados} & RF-8, RF-8.1, RF-8.2 \\
		\textbf{Descripción}          & Permite al usuario personalizar la vista que tiene del conjunto de datos. \\
		\textbf{Precondición}         & El usuario debe haber iniciado sesión y estar en la pantalla de visualizar conjunto. \\
		\textbf{Acciones}             &
		\begin{enumerate}
			\def\labelenumi{\arabic{enumi}.}
			\tightlist
			\item Seleccionar los parámetros deseados para mostrar en el \textit{tooltip} del mapa.
			\item Seleccionar en la tabla el parámetro por el que ordenar los conjuntos de datos.
		\end{enumerate}\\
		\textbf{Postcondición}        & El conjunto de datos se muestra de la forma indicada por el usuario. \\
		\textbf{Excepciones}          & Si el conjunto de datos deja de estar accesible se le redirigirá a la ventana de gestión de errores. \\
		\textbf{Importancia}          & Media \\
		\bottomrule
	\end{tabularx}
	\caption{CU-3.1 Personalizar conjunto.}
\end{table}

% Caso de Uso 4 -> Añadir conjunto.
\begin{table}[p]
	\centering
	\begin{tabularx}{\linewidth}{ p{0.21\columnwidth} p{0.71\columnwidth} }
		\toprule
		\textbf{CU-4}    & \textbf{Añadir conjunto}\\
		\toprule
		\textbf{Versión}              & 1.0    \\
		\textbf{Autor}                & {Alejandro de Castro} \\
		\textbf{Requisitos asociados} & RF-4, RF-4.1 \\
		\textbf{Descripción}          & Permite al usuario añadir un conjunto de datos a la aplicación. \\
		\textbf{Precondición}         & El usuario debe haber iniciado sesión y estar en la pantalla de conjuntos. \\
		\textbf{Acciones}             &
		\begin{enumerate}
			\def\labelenumi{\arabic{enumi}.}
			\tightlist
			\item Rellenar los campos de ciudad, conjunto, enlace, formato y opcionalmente los campos de periodicidad y unidades.
			\item Pulsar el botón de guardar.
		\end{enumerate}\\
		\textbf{Postcondición}        & El conjunto será añadido en la base de datos con los datos introducidos por el usuario. \\
		\textbf{Excepciones}          & Si alguno de los campos obligatorios no está completo, el botón de guardar se encontrará deshabilitado. \\
		\textbf{Importancia}          & Media \\
		\bottomrule
	\end{tabularx}
	\caption{CU-4 Añadir conjunto.}
\end{table}

% Caso de Uso 5 -> Gestionar registros.
\begin{table}[p]
	\centering
	\begin{tabularx}{\linewidth}{ p{0.21\columnwidth} p{0.71\columnwidth} }
		\toprule
		\textbf{CU-5}    & \textbf{Gestionar registros}\\
		\toprule
		\textbf{Versión}              & 1.0    \\
		\textbf{Autor}                & {Alejandro de Castro} \\
		\textbf{Requisitos asociados} & RF-5 \\
		\textbf{Descripción}          & Permite al usuario visualizar una lista de los registros de conjuntos de datos que está guardando el Sistema para él. \\
		\textbf{Precondición}         & El usuario debe haber iniciado sesión. \\
		\textbf{Acciones}             & Pulsar en el botón de registros de la barra de navegación.\\
		\textbf{Postcondición}        & Se mostrará una lista con los registros que el Sistema está realizando para el usuario con las opciones de descargar y eliminar. \\
		\textbf{Excepciones}          & Si no hay ningún registro, se mostrará la lista vacía y la opción de ñadir registros. \\
		\textbf{Importancia}          & Alta \\
		\bottomrule
	\end{tabularx}
	\caption{CU-5 Gestionar registros.}
\end{table}

% Caso de Uso 5.1 -> Añadir registro.
\begin{table}[p]
	\centering
	\begin{tabularx}{\linewidth}{ p{0.21\columnwidth} p{0.71\columnwidth} }
		\toprule
		\textbf{CU-5.1}    & \textbf{Añadir registro}\\
		\toprule
		\textbf{Versión}              & 1.0    \\
		\textbf{Autor}                & {Alejandro de Castro} \\
		\textbf{Requisitos asociados} & RF-5, RF-5.1 \\
		\textbf{Descripción}          & Permite al usuario comenzar el registro de un conjunto de datos determinado con una periodicidad determinada por él. \\
		\textbf{Precondición}         & El usuario debe haber iniciado sesión y debe haber conjuntos de datos almacenados en la base de datos. \\
        \textbf{Acciones}             &
		\begin{enumerate}
			\def\labelenumi{\arabic{enumi}.}
			\tightlist
			\item Seleccionar un conjunto, rellenar la periodicidad y unidades.
			\item Pulsar el botón de guardar.
		\end{enumerate}\\ 
		\textbf{Postcondición}        & Se añadirá la base de datos los datos seleccionados por el usuario y el Sistema comenzará al registro del conjunto de datos con la periodicidad determinada por el usuario. \\
		\textbf{Excepciones}          & Si la periodicidad es menor que la periodicidad de actualización de dicho conjunto se mostrará un mensaje al usuario indicando que debe ser una periodicidad mayor y se desahbilitará el botón de guardar. \\
		\textbf{Importancia}          & Alta \\
		\bottomrule
	\end{tabularx}
	\caption{CU-5.1 Añadir registro.}
\end{table}

% Caso de Uso 5.2 -> Descargar registro.
\begin{table}[p]
	\centering
	\begin{tabularx}{\linewidth}{ p{0.21\columnwidth} p{0.71\columnwidth} }
		\toprule
		\textbf{CU-5.2}    & \textbf{Descargar registro}\\
		\toprule
		\textbf{Versión}              & 1.0    \\
		\textbf{Autor}                & {Alejandro de Castro} \\
		\textbf{Requisitos asociados} & RF-5.2 \\
		\textbf{Descripción}          & Permite al usuario descargar el registro seleccionado. \\
		\textbf{Precondición}         & El usuario debe haber iniciado sesión y debe tener al menos un registro. \\
        \textbf{Acciones}             &
		\begin{enumerate}
			\def\labelenumi{\arabic{enumi}.}
			\tightlist
			\item Seleccionar un registro.
			\item Pulsar el botón de descargar.
		\end{enumerate}\\ 
		\textbf{Postcondición}        & Se descargará en formato \textit{zip} el registro indicado por el usuario. \\
		\textbf{Excepciones}          & Si ocurre algún error durante la descarga el usuario será redirigido a la pantalla de gestión de errores. \\
		\textbf{Importancia}          & Alta \\
		\bottomrule
	\end{tabularx}
	\caption{CU-5.2 Descargar registro.}
\end{table}

% Caso de Uso 5.3 -> Eliminar registro.
\begin{table}[p]
	\centering
	\begin{tabularx}{\linewidth}{ p{0.21\columnwidth} p{0.71\columnwidth} }
		\toprule
		\textbf{CU-5.3}    & \textbf{Eliminar registro}\\
		\toprule
		\textbf{Versión}              & 1.0    \\
		\textbf{Autor}                & {Alejandro de Castro} \\
		\textbf{Requisitos asociados} & RF-5.3 \\
		\textbf{Descripción}          & Permite al usuario dejar de registrar un conjunto de datos. \\
		\textbf{Precondición}         & El usuario debe haber iniciado sesión y debe tener al menos un registro. \\
        \textbf{Acciones}             &
		\begin{enumerate}
			\def\labelenumi{\arabic{enumi}.}
			\tightlist
			\item Seleccionar un registro.
			\item Pulsar el botón de eliminar.
		\end{enumerate}\\ 
		\textbf{Postcondición}        & Se terminará con el registro de ese conjunto para ese usuario. \\
		\textbf{Excepciones}          & Si ocurre algún error durante la eliminación de ese conjunto, el usuario será redirigido a la pantalla de gestión de errores. \\
		\textbf{Importancia}          & Alta \\
		\bottomrule
	\end{tabularx}
	\caption{CU-5.3 Eliminar registro.}
\end{table}

% Caso de Uso 6 -> Reportar error.
\begin{table}[p]
	\centering
	\begin{tabularx}{\linewidth}{ p{0.21\columnwidth} p{0.71\columnwidth} }
		\toprule
		\textbf{CU-6}    & \textbf{Reportar error}\\
		\toprule
		\textbf{Versión}              & 1.0    \\
		\textbf{Autor}                & {Alejandro de Castro} \\
		\textbf{Requisitos asociados} & RF-9 \\
		\textbf{Descripción}          & Permite al usuario dejar constancia de un error producido en la aplicación. \\
		\textbf{Precondición}         & El usuario debe haber iniciado sesión y encontrarse en la pantalla de gestión de errores tras haberse producido un error. \\
        \textbf{Acciones}             &
		\begin{enumerate}
			\def\labelenumi{\arabic{enumi}.}
			\tightlist
			\item Pulsar el botón reportar error.
		\end{enumerate}\\ 
		\textbf{Postcondición}        & Se guardará el error encontrado, la fecha y el usuario para su revisión y el usuario. \\
		\textbf{Excepciones}          & No hay excepciones en esta página al ser crítica no se permiten fallos. \\
		\textbf{Importancia}          & Alta \\
		\bottomrule
	\end{tabularx}
	\caption{CU-6 Reportar error.}
\end{table}

% Caso de Uso 7 -> Gestionar favoritos.
\begin{table}[p]
	\centering
	\begin{tabularx}{\linewidth}{ p{0.21\columnwidth} p{0.71\columnwidth} }
		\toprule
		\textbf{CU-7}    & \textbf{Gestionar favoritos}\\
		\toprule
		\textbf{Versión}              & 1.0    \\
		\textbf{Autor}                & {Alejandro de Castro} \\
		\textbf{Requisitos asociados} & RF-3, RF-3.1 \\
		\textbf{Descripción}          & Permite al usuario visualizar una lista de los conjuntos que tiene en favoritos. \\
		\textbf{Precondición}         & El usuario debe haber iniciado sesión. \\
        \textbf{Acciones}             &
		\begin{enumerate}
			\def\labelenumi{\arabic{enumi}.}
			\tightlist
			\item Pulsar en el botón \textit{home} de la barra de navegación.
		\end{enumerate}\\ 
		\textbf{Postcondición}        & Se mostrará una lista con los conjuntos establecidos como favoritos por el usuario. \\
		\textbf{Excepciones}          & En caso de no tener ningún conjunto marcado como favoritos se mostrará una lista vacía. \\
		\textbf{Importancia}          & Alta \\
		\bottomrule
	\end{tabularx}
	\caption{CU-7 Gestionar favoritos.}
\end{table}

% Caso de Uso 7.1 -> Marcar favorito.
\begin{table}[p]
	\centering
	\begin{tabularx}{\linewidth}{ p{0.21\columnwidth} p{0.71\columnwidth} }
		\toprule
		\textbf{CU-7.1}    & \textbf{Marcar favorito}\\
		\toprule
		\textbf{Versión}              & 1.0    \\
		\textbf{Autor}                & {Alejandro de Castro} \\
		\textbf{Requisitos asociados} & RF-3 \\
		\textbf{Descripción}          & Permite al usuario marcar como favorito un conjunto de datos. \\
		\textbf{Precondición}         & \begin{itemize}
		    \item El usuario debe haber iniciado sesión. 
            \item El usuario debe encontrarse en la pantalla de visualización de un conjunto. 
            \item El usuario no debe tener ese conjunto marcado como favorito previamente. 
		\end{itemize} \\
        \textbf{Acciones}             &
		\begin{enumerate}
			\def\labelenumi{\arabic{enumi}.}
			\tightlist
			\item Pulsar en el botón en forma de corazón negro.
		\end{enumerate} \\ 
		\textbf{Postcondición}        & Se añadirá el conjunto a la lista de favoritos de ese usuario en la base de datos y el botón en forma de corazón negro quedará coloreado de rojo. \\
		\textbf{Excepciones}          & En caso de no poderse añadir a favoritos se direccionara al usuario a la pantalla de gestión de errores. \\
		\textbf{Importancia}          & Media \\
		\bottomrule
	\end{tabularx}
	\caption{CU-7.1 Marcar favorito.}
\end{table}

% Caso de Uso 7.2 -> Desmarcar favorito.
\begin{table}[p]
	\centering
	\begin{tabularx}{\linewidth}{ p{0.21\columnwidth} p{0.71\columnwidth} }
		\toprule
		\textbf{CU-7.2}    & \textbf{Desmarcar favorito}\\
		\toprule
		\textbf{Versión}              & 1.0    \\
		\textbf{Autor}                & {Alejandro de Castro} \\
		\textbf{Requisitos asociados} & RF-3 \\
		\textbf{Descripción}          & Permite al usuario desmarcar como favorito un conjunto de datos. \\
		\textbf{Precondición}         & \begin{itemize}
		    \item El usuario debe haber iniciado sesión. 
            \item El usuario debe encontrarse en la pantalla de visualización de un conjunto. 
            \item El usuario debe tener ese conjunto marcado como favorito previamente. 
		\end{itemize} \\
        \textbf{Acciones}             &
		\begin{enumerate}
			\def\labelenumi{\arabic{enumi}.}
			\tightlist
			\item Pulsar en el botón en forma de corazón rojo.
		\end{enumerate}\\ 
		\textbf{Postcondición}        & Se eliminará el conjunto a la lista de favoritos de ese usuario en la base de datos y el botón en forma de corazón cambiará de color a negro. \\
		\textbf{Excepciones}          & En caso de no poderse eliminar de favoritos se direccionará al usuario a la pantalla de gestión de errores. \\
		\textbf{Importancia}          & Media \\
		\bottomrule
	\end{tabularx}
	\caption{CU-7.2 Desmarcar favorito.}
\end{table}

% Caso de Uso 8 -> Visualizar ayuda.
\begin{table}[p]
	\centering
	\begin{tabularx}{\linewidth}{ p{0.21\columnwidth} p{0.71\columnwidth} }
		\toprule
		\textbf{CU-8}    & \textbf{Visualizar ayuda}\\
		\toprule
		\textbf{Versión}              & 1.0    \\
		\textbf{Autor}                & {Alejandro de Castro} \\
		\textbf{Requisitos asociados} & RF-10 \\
		\textbf{Descripción}          & Permite al usuario visualizar una pantalla con información sobre el funcionamiento de la aplicación y el contacto de los administradores principales. \\
		\textbf{Precondición}         & El usuario debe haber iniciado sesión. \\
        \textbf{Acciones}             &
		\begin{enumerate}
			\def\labelenumi{\arabic{enumi}.}
			\tightlist
			\item Pulsar en el botón ayuda de la barra de navegación.
		\end{enumerate}\\ 
		\textbf{Postcondición}        & Se mostrará una pantalla con información sobre el funcionamiento de la aplicación y el contacto de los administradores principales. \\
		\textbf{Excepciones}          & Este apartado es crítico, por lo que no hay excepciones permitidas. \\
		\textbf{Importancia}          & Alta \\
		\bottomrule
	\end{tabularx}
	\caption{CU-8 Visualizar ayuda.}
\end{table}
\apendice{Especificación de diseño}

\section{Introducción}
En este anexo se define la resolución de los objetivos y especificaciones de los anexos anteriores.

Su estructura es la siguiente:
\begin{itemize}
    \item \textbf{Diseño de datos:} se detalla el modelado de datos de la aplicación mediante el diagrama Entidad-Relación \ref{e-r}, el modelo relacional \ref{modelo-relacional} y el diccionario de datos \ref{diccionario-de-datos} de la aplicación.
    \item \textbf{Diseño procedimental:} se detalla mediante diagramas de secuencia las distintas funcionalidades de la aplicación.
    \item \textbf{Diseño arquitectónico:} se muestra un diagrama de la arquitectura de la aplicación.
    \item \textbf{Diseño gráfico:} se muestra el diagrama de interfaces desarrollado como guía para el producto final.
\end{itemize}

\section{Diseño de datos}
\subsection{Diagrama Entidad-Relación}\label{e-r}

Un diagrama entidad-relación (ERD) describe los datos o información de un sistema mediante modelado de sus entidades, atributos y las relaciones entre esas entidades. \cite{Chen1976}

Como se puede ver en el diagrama \ref{fig:E_R}, existen 4 entidades:
\begin{itemize}
    \item \textbf{Usuarios:} poseen identificador, usuario, contraseña, nombre completo, favoritos y rol.
    \item \textbf{Datos:} poseen identificador, ciudad, característica, enlace, formato y periodicidad.
    \item \textbf{Registros:} poseen identificador, ciudad, característica, formato, periodicidad y \textit{pid}.
    \item \textbf{Traducciones:} poseen identificador, usuario, contraseña, nombre completo, favoritos y rol.
\end{itemize}

Usuarios se relaciona con Datos con \textit{id usuario}, guardado en Datos. Un usuario puede añadir ninguno o N conjuntos de datos, por lo tanto, cada conjunto de datos almacenado debe tener como mínimo y máximo un \textit{id usuario} para identificar quién guardó ese conjunto.

Usuarios se relaciona con Registros con \textit{id usuario}, guardado en Registros. Un usuario puede tener ninguno o N registros para él, por lo tanto, cada registro almacenado debe tener como mínimo y máximo un \textit{id usuario} para identificar para quién es dicho registro.

Traducciones no tiene relación con ninguna tabla, ya que solo pueden escribir administradores en esa tabla, por lo tanto, no es interesante guardar el \textit{id usuario} que añadió la traducción en cuestión, ya que los administradores son un número muy reducido de usuarios y de confianza. Y el número de traducciones es muy elevado por lo que se guardaría información repetida e irrelevante de manera muy recurrente.

\imagen{E_R}{Diagrama Entidad-Relación}


\subsection{Modelo Relacional}\label{modelo-relacional}

Un modelo relacional es una forma de estructurar y consultar datos en una base de datos. En este modelo, los datos se organizan en tablas (o relaciones), donde cada tabla está compuesta por filas (tuplas) y columnas (atributos). \cite{Codd1970}

El modelo relacional de esta aplicación lo podemos observar en la imagen \ref{fig:ModeloRelacional}

\imagen{ModeloRelacional}{Modelo Relacional}


\subsection{Diccionario de datos}\label{diccionario-de-datos}

Un diccionario de datos es una herramienta crucial en la gestión de bases de datos que almacena definiciones y descripciones de los elementos de datos utilizados en el sistema. Proporciona una visión centralizada y coherente de los datos, incluyendo detalles sobre los nombres, tipos, formatos, y relaciones entre los elementos de datos. \cite{Kent1983}

La tabla~\ref{tabla-dic-usuarios} contiene la información de la tabla de Usuarios.

\begin{table}
	\scalebox{0.80}{
		\begin{tabular}{@{}p{10em} p{6em} p{6em} p{20em}@{}}
			\toprule
			\textbf{Nombre} & \textbf{Tipo} & \textbf{Columna} & \textbf{Descripción}\\
			\midrule
			\texttt{\textbf{\underline{id}}} & \texttt{SMALLINT(3)} & \texttt{\textbf{\underline{PK}}} & Identificador. Generado automáticamente. \\
			\texttt{usuario} &  \texttt{VARCHAR(20)} & \texttt{NOT NULL} & Nombre de usuario proporcionado en el registro.\\
			\texttt{contraseña} &  \texttt{CHAR(255)} & \texttt{NOT NULL} & Contraseña del usuario cifrada.\\
			\texttt{nombre\_completo} & \texttt{VARCHAR(50)} & \texttt{NOT NULL} & Nombre completo del usuario.\\
			\texttt{favoritos} & \texttt{text} & \texttt{NOT NULL} & Lista de conjuntos favoritos del usuario separada por comas. \\
			\texttt{rol} & \texttt{VARCHAR(15)} & \texttt{NOT NULL} & Rol del usuario. \\
			\bottomrule
		\end{tabular}
	}
	\caption{Diccionario de Datos - Tabla Usuarios.}
    \label{tabla-dic-usuarios}
\end{table}

La tabla~\ref{tabla-dic-datos} contiene la información de la tabla de Datos.
\begin{table}
	\scalebox{0.80}{
		\begin{tabular}{@{}p{10em} p{6em} p{6em} p{20em}@{}}
			\toprule
			\textbf{Nombre} & \textbf{Tipo} & \textbf{Columna} & \textbf{Descripción}\\
			\midrule
			\texttt{\textbf{\underline{id}}} & \texttt{SMALLINT(6)} & \texttt{\textbf{\underline{PK}}} & Identificador. Generado automáticamente. \\
			\texttt{ciudad} &  \texttt{VARCHAR(25)} & \texttt{NOT NULL} & Ciudad en la que se encuentra el conjunto de datos.\\
			\texttt{característica} &  \texttt{VARCHAR(100)} & \texttt{NOT NULL} & Nombre del conjunto de datos.\\
			\texttt{enlace} &  \texttt{VARCHAR(10000)} & \texttt{NOT NULL} & Enlace al conjunto de datos.\\
			\texttt{id\_usuario} & \texttt{SMALLINT(3)} & \texttt{NOT NULL} & Identificador del usuario que ha añadido el conjunto.\\
			\texttt{formato} & \texttt{VARCHAR(8)} & \texttt{NOT NULL} & Formato del conjunto de datos. \\
			\texttt{periodicidad} & \texttt{INTEGER(250)} & \texttt{NOT NULL} & Periodicidad de actualización del conjunto de datos en segundos. \\
			\bottomrule
		\end{tabular}
	}
	\caption{Diccionario de Datos - Tabla Datos.}
    \label{tabla-dic-datos}
\end{table}

La tabla~\ref{tabla-dic-registros} contiene la información de la tabla de Registros.
\begin{table}
	\scalebox{0.80}{
		\begin{tabular}{@{}p{10em} p{6em} p{6em} p{20em}@{}}
			\toprule
			\textbf{Nombre} & \textbf{Tipo} & \textbf{Columna} & \textbf{Descripción}\\
			\midrule
			\texttt{\textbf{\underline{id}}} & \texttt{SMALLINT(3)} & \texttt{\textbf{\underline{PK}}} & Identificador. Generado automáticamente. \\
			\texttt{ciudad} &  \texttt{VARCHAR(25)} & \texttt{NOT NULL} & Ciudad en la que se encuentra el conjunto de datos.\\
			\texttt{característica} &  \texttt{VARCHAR(100)} & \texttt{NOT NULL} & Nombre del conjunto de datos.\\
			\texttt{id\_usuario} & \texttt{SMALLINT(3)} & \texttt{NOT NULL} & Identificador del usuario al que pertenece el registro.\\
			\texttt{formato} & \texttt{VARCHAR(8)} & \texttt{NOT NULL} & Formato del conjunto de datos. \\
			\texttt{periodicidad} & \texttt{INTEGER(250)} & \texttt{NOT NULL} & Periodicidad de captura de datos. \\
            \texttt{pid} &  \texttt{INTEGER(6)} & \texttt{NOT NULL} & Identificación del proceso que realiza captura de datos.\\
			\bottomrule
		\end{tabular}
	}
	\caption{Diccionario de Datos - Tabla Registros.}
    \label{tabla-dic-registros}
\end{table}

La tabla~\ref{tabla-dic-traducciones} contiene la información de la tabla de Traducciones.
\begin{table}
	\scalebox{0.80}{
		\begin{tabular}{@{}p{10em} p{6em} p{6em} p{20em}@{}}
			\toprule
			\textbf{Nombre} & \textbf{Tipo} & \textbf{Columna} & \textbf{Descripción}\\
			\midrule
			\texttt{\textbf{\underline{id}}} & \texttt{SMALLINT(3)} & \texttt{\textbf{\underline{PK}}} & Identificador. Generado automáticamente. \\
			\texttt{original} &  \texttt{VARCHAR(100)} & \texttt{NOT NULL} & Palabra o frase a traducir.\\
            \texttt{traducción} &  \texttt{VARCHAR(100)} & \texttt{NOT NULL} & Palabra o frase traducida.\\
			\bottomrule
		\end{tabular}
	}
	\caption{Diccionario de Datos - Tabla Traducciones.}
    \label{tabla-dic-traducciones}
\end{table}

\section{Diseño procedimental}

Un diagrama de secuencias es un tipo de diagrama de interacción. Este diagrama se utiliza para representar cómo los objetos interactúan en un proceso específico a lo largo del tiempo, mostrando la secuencia de mensajes intercambiados entre los objetos para llevar a cabo una funcionalidad particular. \cite{Booch1998}

Gracias al diagrama de secuencias se puede describir en detalle cómo se lleva a cabo una funcionalidad específica. Es por ello, por lo que se han elaborado los siguientes diagramas de secuencia, los cuales han servido de guía para el desarrollo de la aplicación.

\subsection{Crear cuenta}
\imagen{UC-1}{Procedimiento de crear una cuenta}

\subsection{Iniciar sesión}
\imagen{UC-2}{Procedimiento de iniciar sesión}

\subsection{Visualizar conjunto}
\imagen{UC-3}{Procedimiento de visualizar un conjunto}

\subsection{Añadir conjunto}
\imagen{UC-4}{Procedimiento de añadir un conjunto}

\subsection{Mostrar registros}
\imagen{UC-5}{Procedimiento de mostrar los registros}

\subsection{Añadir registro}
\imagen{UC-5.1}{Procedimiento de añadir un registro}

\subsection{Descargar registro}
\imagen{UC-5.2}{Procedimiento de descargar registro}

\subsection{Eliminar registro}
\imagen{UC-5.3}{Procedimiento de eliminar registro}

\subsection{Eliminar conjunto}
\imagen{UC-10}{Procedimiento de eliminar un conjunto}

\subsection{Añadir traducciones}
\imagen{UC-11.1}{Procedimiento de añadir una traducción}


\section{Diseño arquitectónico}

\subsection{Estructura del \textit{software} del proyecto}
\imagen{Estructura}{Estructura del proyecto}

\subsection{Arquitectura de la aplicación}
\imagen{Diseño arquitectónico}{Diseño arquitectónico de la aplicación}

\section{Diseño gráfico}

Se ha elaborado un

\subsection{Crear cuenta}
\imagen{crear_cuenta_pantalla}{Crear cuenta - Pantalla}

\subsection{Login}
\imagen{login_pantalla}{Login - Pantalla}

\subsection{Home - Usuario}
\imagen{home_pantalla}{Home usuario - Pantalla}

\subsection{Home - Administrador}
\imagen{home_admin_pantalla}{Home administrador - Pantalla}

\subsection{Conjuntos}
\imagen{conjuntos_pantalla}{Conjuntos - Pantalla}

\subsection{Registros}
\imagen{registros_pantalla}{Registros - Pantalla}

\subsection{Conjunto}
\imagen{conjunto_pantalla}{Conjunto - Pantalla}

\subsection{Ayuda}
\imagen{ayuda_pantalla}{Ayuda - Pantalla}

\apendice{Documentación técnica de programación}

\section{Introducción}

\section{Estructura de directorios}

El repositorio de \textit{GitHub}\footnote{\url{https://github.com/AlejandroDeCastro/TFG}} donde se en encuentra el proyecto, tiene la siguiente estructura:
\begin{itemize}
    \item \textbf{CODE:} Directorio donde se encuentra el código de la aplicación. Tiene la siguiente estructura:
    \begin{itemize}
        \item \textbf{assets:} Contiene los estilos \textit{css} de \textit{Dash}:
        \begin{itemize}
            \item \textbf{stylesheet.css:} Fichero \textit{css} que contiene el diseño de las vistas hechas con \textit{Dash}.
            \item \textbf{s1.css:} Fichero \textit{css} que contiene el diseño de las gráficas de las vistas hechas con \textit{Dash}.
        \end{itemize}
        \item \textbf{models:} Contiene los modelos de las entidades que ejecutan el programa y las acciones relacionadas con ellas:
        \begin{itemize}
            \item \textbf{entidades:} Directorio que contiene las distintas entidades del sistema:
            \begin{itemize}
                \item \textbf{Usuario.py:} Contiene la clase "Usuario". 
            \end{itemize}
            \item \textbf{ModeloUsuario.py:} Contiene la clase "ModeloUsuario" en \textit{Python}, la cual alberga todas acciones relacionadas con el usuario.
        \end{itemize}
        \item \textbf{static:} Directorio donde se encuentran los ficheros que dan soporte al \textit{front-end} de la aplicación y algunas funcionalidades en las propias pantallas de la aplicación:
        \begin{itemize}
            \item \textbf{css:} Directorio donde se encuentran todos los ficheros \textit{css} de la aplicación, encargados de aportar la estética de las distintas vistas.
            \item \textbf{img:} Directorio con las distintas imágenes de la aplicación.
            \item \textbf{js:} Directorio que contiene todos los ficheros \textit{js} de la aplicación, son los encargados de añadir algunas funcionalidades a las distintas vistas.
        \end{itemize}
        \item \textbf{templates:} Directorio que contiene los ficheros \textit{html} de la aplicación, estos ficheros aportan las vistas de la aplicación.
        \item \textbf{app.py:} \textit{Script} principal de la aplicación, programado en \textit{Python}, contiene todas las rutas de la aplicación.
        \item \textbf{database.py:} Fichero de \textit{Python} que contiene el conector con la base de datos.
        \item \textbf{gestor.py:} Fichero de \textit{Python} que contiene todas las interacciones con la base de datos, menos las que están relacionadas con el usuario. También es el \textit{script} encargado de la gestión de registros.
    \end{itemize}
    \item \textbf{Documentación:} Directorio donde se encuentra toda la documentación del proyecto, organizada en documentos \LaTeX{}
    \item \textbf{Registros:} Directorio donde se almacenan los registros solicitados por cada usuario.
    \item \textbf{Server:} Directorio donde se encuentra el \textit{OCB} y todos los ficheros relacionados con él:
    \begin{itemize}
        \item \textbf{Server/Ficheros:} Directorio donde se encuentra un histórico de los ficheros \textit{JSON} del simulador.
        \item \textbf{menu.bat:} \textit{Script} que proporciona al usuario un menú desde el cual gestionar el \textit{OCB} en \textit{Windwos}.
        \item \textbf{menu.sh:} \textit{Script} que proporciona al usuario un menú desde el cual gestionar el \textit{OCB} en \textit{Linux}.
        \item \textbf{generadorDatos.py:} Simulador que proporciona datos al \textit{OCB}.
        \item \textbf{limpiador.py:} \textit{Script} que eliminar todas las entidades almacenadas en el \textit{OCB}.
        \item \textbf{simulador.txt:} Fichero de texto plano que almacena el \textit{id} del proceso que está ejecutando el simulador para detenerlo cuando se ejecuta en segundo plano.
        \item \textbf{transformadorDatos.py:} \textit{Script}que transforma ficheros \textit{JSON} al modelo \textit{NGSI}, actualmente en desuso.
    \end{itemize}
    \item \textbf{db:} Directorio que contiene la base de datos de la aplicación.
    \begin{itemize}
        \item \textbf{tfg.sql:} Base de datos de la aplicación.
    \end{itemize}
    \item \textbf{Dockerfile:} Fichero de texto plano que contiene una serie de instrucciones necesarias para crear la imagen \textit{Docker} que contiene la aplicación, para su distribución o despliegue en un servidor real.
    \item \textbf{LICENSE:} Licencia para la  distribución del código de la aplicación.
    \item \textbf{Readme.md:} Fichero que proporciona información útil sobre el proyecto, como la licencia, enlaces de interés, descripción, etc.
    \item \textbf{docker-compose.yml:} Fichero de configuración para la creación del \textit{Docker} mencionado anteriormente.
    \item \textbf{error\_log.txt:} Fichero de texto plano donde se almacenan los reportes de los usuarios cuando ha ocurrido un error.
    \item \textbf{requirements.txt:} Bibliotecas esenciales y sus versiones para la ejecución de la aplicación.
\end{itemize}

Se han omitido los directorios \textbf{\_\_pycache\_\_} que contienen \textit{scripts} ya compilados de \textit{Python}, se genera automáticamente. No son relevantes.


\section{Manual del programador}

\section{Compilación, instalación y ejecución del proyecto}

\section{Pruebas del sistema}

\apendice{Documentación de usuario}

\section{Introducción}
En este anexo


\section{Requisitos de usuarios}

Navegador

\section{Instalación}

La aplicación al ser desplegada carece de 
El usuario no tiene que descargar

\section{Manual del usuario}


\subsection{Iniciar sesión}\label{inicio-sesión}
Cuando el usuario entra en el enlace de la aplicación, lo primero que se le muestra es la pantalla de inicio \ref{fig:Pantalla - Login} de sesión donde debe introducir su usuario y contraseña.
\imagen{Pantalla - Login}{Pantalla de inicio de sesión}

En caso de introducir mal la contraseña se le notificará al usuario con un mensaje de error como se muestra en la imagen \ref{fig:Pantalla - Login - Contraseña incorrecta}.
\imagen{Pantalla - Login - Contraseña incorrecta}{Mensaje de contraseña incorrecta}

Y en caso de no existir un usuario con ese nombre en la base de datos se le notificará también con un error como el que se muestra en la imagen \ref{fig:Pantalla - Login - Usuario no encontrado}.
\imagen{Pantalla - Login - Usuario no encontrado}{Mensaje de usuario no encontrado}


\subsection{Crear una cuenta}\label{crear-cuenta}

Si el usuario no tiene una cuenta en la aplicación puede crear una pulsando en el enlace que dice "¿No tiene cuenta aún? Créala aquí", será redirigido a la pantalla  \ref{fig:Pantalla - Crear cuenta}, donde puede crear una cuenta introduciendo un nombre de usuario, nombre completo y la contraseña dos veces.
\imagen{Pantalla - Crear cuenta}{Pantalla de registro}

En caso de existir un usuario con ese nombre de usuario se le notificará con un mensaje como el de la imagen \ref{fig:Pantalla - Crear cuenta - Usuario existe}.
\imagen{Pantalla - Crear cuenta - Usuario existe}{Mensaje usuario ya registrado}

Y en el caso de no coincidir las dos contraseñas también se le notificará al usuario.
\imagen{Pantalla - Crear cuenta - Contraseñas no coinciden}{Mensaje contraseñas no coinciden}

Una vez introducido los datos correctamente, el usuario solo tiene que pulsar el botón "Registrar" y la cuenta será creada. Posteriormente será redirigido la pantalla de login \ref{fig:Pantalla - Login} donde deberá introducir sus datos para continuar.


\subsection{Pantalla principal}\label{home}
Una vez iniciada sesión en la aplicación el usuario será redirigido a la pantalla principal, la cual varía según el rol del usuario. Si se acaba de crear la cuenta por defecto será un usuario estándar y la pantalla que verá será la representada en la imagen \ref{fig:Pantalla - Home}:
\imagen{Pantalla - Home}{Pantalla principal para un usuario estándar}

Por el contrario, si es un usuario con permiso de administrador, esta será la pantalla principal para ese usuario:
\imagen{Pantalla - Home - Administrador}{Pantalla principal para un usuario administrador}

En esta pantalla principal ambos usuarios pueden ver en la parte inferior izquierda los conjuntos que tienen marcados en favoritos y acceder a ellos pulsando sobre ellos o seleccionando la ciudad y el conjunto en el cuadrado inferior de la parte derecha de la pantalla. Ambas opciones redirigirán al usuario a la pantalla de visualizar conjunto \ref{visualizar-conjunto}.

Los usuarios estándar podrán consultar dentro del recuadro verde que como se muestra en la imagen \ref{fig:Pantalla - Home - Registros} los registros que tienen en ese momento guardando.
\imagen{Pantalla - Home - Registros}{Cuadro de registros de la pantalla principal}

Los usuarios con rol de administrador tendrán en su defecto un menú con tres opciones que redirigirán al usuario a la interfaz de gestión de usuarios \ref{gestión-usuarios}, la lista de conjuntos disponibles \ref{consultar-conjuntos} o la pantalla de gestión de traducciones \ref{gestión-traducciones}, como se muestra en la imagen \ref{fig:Pantalla - Home - Administrador - Opciones}.
\imagen{Pantalla - Home - Administrador - Opciones}{Opciones del administrador en la pantalla principal}

Ambos tipos de usuarios, en la barra de navegación, podrán consultar conjuntos disponibles \ref{consultar-conjuntos} pulsando en el botón "Conjuntos", consultar los registros que tienen actualmente guardando, pulsando en el botón "Registros" \ref{consultar-registros}, consultar la pantalla de ayuda \ref{consulta-ayuda} con el botón "Ayuda" o cerrar la sesión pulsando el botón en rojo "Cerrar sesión". Si el usuario selecciona esta opción será redirigido a la pantalla de inicio de sesión \ref{inicio-sesión}

\subsection{Consultar los conjuntos de datos disponibles}\label{consultar-conjuntos}
Desde la pantalla \ref{fig:Pantalla - Conjuntos} el usuario puede consultar una lista con los conjuntos disponibles en la aplicación y acceder a ellos pulsando en el botón "consultar" azul de cada conjunto. Será direccionado al la visualización del conjunto seleccionado \ref{visualizar-conjunto}.
\imagen{Pantalla - Conjuntos}{Pantalla que mmuestra al usuario todos los conjuntos disponibles}

También el usuario puede ordenar de manera ascendente o descendente los conjuntos en función de la ciudad, el tipo de conjunto o el formato, pulsando en el nombre del campo en el cabecero de la tabla.

Si desea añadir un conjunto el usuario solo tiene que rellenar los campos de la parte superior de  la pantalla y posteriormente pulsar el botón "Añadir", y automáticamente se añadirá el conjunto a la lista.

Los administradores tienen un botón rojo "Eliminar" en cada conjunto. Esto permite a los administradores eliminar el conjunto que seleccionen, podemos observar la pantalla de los administradores en \ref{fig:Pantalla - Conjuntos - Administrador}.
\imagen{Pantalla - Conjuntos - Administrador}{Pantalla de conjuntos de un administrador}

\subsection{Consultar registros}\label{consultar-registros}
Desde la pantalla de visualizar registros \ref{fig:Pantalla - Registros}, se pueden ver todos los registros que ese usuario tiene en ese momento.
También, se pueden añadir nuevos registros seleccionando el conjunto, añadiendo una periodicidad y pulsando en el botón "Guardar". Automáticamente se comenzará con la captura de datos.
Por último también se pueden eliminar los registros que el usuario quiera dejar de capturar pulsando el botón "Eliminar"

\subsection{Visualizar conjunto en tiempo real}\label{visualizar-conjunto}



\subsection{Gestionar traducciones}\label{gestión-traducciones}
Un administrador desde la pantalla \ref{fig:Pantalla - Traducciones} puede consultar todas las traducciones que existen en la aplicación. Estas traducciones son aplicadas a los campos de los conjuntos de datos.
El administrador puede añadir nuevas traducciones rellenando los campos "Campo original" y "Traducción" y pulsando el botón "Añadir".
\imagen{Pantalla - Traducciones}{Pantalla de gestión de traducciones}

\subsection{Gestionar usuarios}\label{gestión-usuarios}
Desde la pantalla de gestión de usuarios \ref{fig:Pantalla - Traducciones} un administrador puede consultar los usuarios registrados en la aplicación, modificar su rol mediante el desplegable del campo "Rol" o eliminarlo desde el botón de "Eliminar".
\imagen{Pantalla - Usuarios}{Pantalla de gestión de usuarios}

\subsection{Ayuda}\label{consulta-ayuda}
Gracias a la pantalla de ayuda \ref{fig:Pantalla - Ayuda}, se pueden consultar las preguntas más frecuentes realizadas por otros usuarios y sus respuestas. También se puede acceder en el apartado de enlaces de interés a una guía elaborada con la herramienta \textit{GitBook}\footnote{\url{https://www.gitbook.com}}, al repositorio del código de la aplicación en \textit{GitBook}\footnote{\url{https://www.github.com}} o abrir un correo nuevo con un administrador como destinatario para contactar con él.
\imagen{Pantalla - Ayuda}{Pantalla de ayuda}

\subsection{Errores}\label{errores}
Si en algún punto de la aplicación algo no saliese como debería, se mostraría la pantalla de error \ref{fig:Pantalla - Error}, en la cuál se puede observar cual ha sido el error y reportarlo si se desea, también el usuario puede volver directamente la pantalla principal \ref{fig:Pantalla - Home}
\imagen{Pantalla - Error}{Pantalla de error}
\apendice{Anexo de sostenibilización curricular}

\begin{comment}
Este anexo incluirá una reflexión personal del alumnado sobre los aspectos de la sostenibilidad que se abordan en el trabajo.
Se pueden incluir tantas subsecciones como sean necesarias con la intención de explicar las competencias de sostenibilidad adquiridas durante el alumnado y aplicadas al Trabajo de Fin de Grado.

Más información en el documento de la CRUE \url{https://www.crue.org/wp-content/uploads/2020/02/Directrices_Sosteniblidad_Crue2012.pdf}.

Este anexo tendrá una extensión comprendida entre 600 y 800 palabras.
\end{comment}

\section{Introducción}
En este anexo se incluye una reflexión personal sobre los aspectos de la sostenibilidad que se abordan en el trabajo.

La sostenibilidad es un aspecto fundamental en el desarrollo de aplicaciones tecnológicas, especialmente aquellas que manejan datos abiertos. Este capítulo describe cómo la aplicación desarrollada durante este proyecto contribuye a la sostenibilidad ambiental, social y económica.

\section{Sostenibilidad Ambiental}
\subsection{Impacto Positivo}
La aplicación facilita el acceso a datos sobre la calidad del aire, el consumo de energía, y otros indicadores ambientales que estén disponibles. Esta información permite a las autoridades y a los ciudadanos tomar decisiones informadas que promuevan un entorno urbano más saludable.

\subsection{Gestión de Recursos}
Al proporcionar datos abiertos, la aplicación ayuda a mejorar la planificación urbana y la eficiencia en el uso de recursos, contribuyendo así a una gestión más sostenible. Como restricciones de tráfico en determinados distritos, reforzar líneas de transporte cuyos aforos se aproximen al máximo o incluso crear nuevas líneas de transporte público en aquellas zonas donde sea necesario.

\section{Sostenibilidad Social}
\subsection{Transparencia y Participación}
El acceso libre a datos públicos fomenta la transparencia gubernamental y la participación ciudadana, empoderando a los individuos para involucrarse en la toma de decisiones que afectan a su comunidad.

\subsection{Inclusión}
La aplicación promueve la equidad al proporcionar acceso igualitario a información relevante, reduciendo así las desigualdades sociales.

\section{Sostenibilidad Económica}
\subsection{Eficiencia y Ahorro}
El uso de datos abiertos permite optimizar los recursos económicos de la ciudad, mejorando la eficiencia en diferentes sectores y fomentando un desarrollo económico sostenible.

\section{Educación y Conciencia}
\subsection{Potencial Educativo}
La aplicación actúa como una herramienta educativa, ofreciendo datos que pueden ser utilizados en investigaciones académicas y estudios.

\subsection{Sensibilización}
La disponibilidad de datos abiertos aumenta la conciencia sobre problemas ambientales y sociales, incentivando a la ciudadanía a adoptar prácticas más sostenibles.

\section{Innovación y Tecnología}
\subsection{Uso de Tecnologías Sostenibles}
El desarrollo de la aplicación incluye la implementación de tecnologías eficientes y el uso de servidores de bajo consumo energético, asegurando un impacto ambiental reducido.

\section{Conclusiones}
El proyecto aborda competencias esenciales como la comprensión del impacto de la actividad profesional en el medio ambiente y la sociedad, el trabajo en equipos multidisciplinares, y la aplicación de un enfoque holístico para resolver problemas socioambientales.

Esto se puede ver reflejado en cada uno de los distintos conjuntos de datos que se pueden incluir en la aplicación, ya que, todos ellos contribuyen a la sostenibilidad, ya sea económica, social o ambiental. 




\bibliographystyle{plain}
\bibliography{bibliografiaAnexos}

\end{document}
